% Reversibility without Injectivity
% Springer LNCS -- contributed paper, 15 pages of body, references uncounted.
%
% Build:  pdflatex paper && bibtex paper && pdflatex paper && pdflatex paper
%
% llncs.cls (v2.26) and splncs04.bst are alongside this file, from CTAN.
% `envcountsame` puts every result on one counter, which is what the text assumes.

\documentclass[runningheads,envcountsame]{llncs}

\usepackage[T1]{fontenc}
\usepackage[utf8]{inputenc}
\usepackage{amsmath}
\usepackage{amssymb}
\usepackage{listings}

% amsmath warns "Unable to redefine math accent \vec" under llncs.  It is
% benign: llncs deliberately defines \vec as bold rather than as an arrow
% accent, and amsmath cannot override that.  Nothing here uses \vec -- the
% n-tuple of empty relations is written out -- so the class's definition stands.

% LNCS \textwidth is 122mm.  At \footnotesize, cmtt gives 4pt per character,
% so roughly 86 characters fit; no line below exceeds 80.  `xleftmargin' is
% deliberately absent -- it costs width the listings do not have to spare.
\lstset{basicstyle=\ttfamily\footnotesize,columns=fullflexible,keepspaces=true,
        showstringspaces=false,frame=none,aboveskip=0.6em,belowskip=0.6em}

% Denotation brackets without depending on stmaryrd.
\newcommand{\den}[1]{\ensuremath{[\![#1]\!]}}
\newcommand{\iso}{\leftrightarrow}
\newcommand{\conv}{^{\circ}}
\newcommand{\Rel}{\textbf{Rel}}
\newcommand{\PInj}{\textbf{PInj}}
% An inference rule, laid out without mathpartir.
\newcommand{\rul}[2]{\ensuremath{\dfrac{\;#1\;}{\;#2\;}}}
% closure operator, written `^' in the concrete syntax
\newcommand{\clo}{\ensuremath{^{\wedge}}}

\begin{document}

\title{Reversibility without Injectivity}
\subtitle{A Point-Free Reversible Language Interpreted in \textbf{Rel}}
\titlerunning{Reversibility without Injectivity}

% [AUTHORS] The 1993 language this one descends from is credited in its thesis to
% Joep Rous and Paul Jansen, and the paper listed there as forthcoming was to be
% joint.  Whether this paper is single- or joint-authored is not settled here.
\author{Joep Rous\inst{1}}
\authorrunning{J. Rous}
\institute{Affiliation \email{ubrowz@gmail.com}}

\maketitle

\begin{abstract}
Reversible functional languages require the clauses of a program to be
non-overlapping, so that backward execution is deterministic. Inv assumes the
condition and leaves checking it as future work, Theseus makes it the single rule
a programmer must maintain, RFUN relaxes it into a first-match ordering, PisoLang
decides it by unification, and Chardonnet et al.\ formalise it as an orthogonality
premise in the typing rule. We ask what happens if it is simply not imposed.

The answer is 42, a point-free language interpreted in the category \Rel{} of sets
and relations, whose defining law is $x \in P(y) \iff y \in \mathrm{inv}(P)(x)$.
Giving up injectivity makes the metatheory \emph{smaller}: inversion becomes a
total, involutive operation with no side conditions, no well-formedness checks and
no proof obligations, computed by an eight-line structural recursion and
eliminated at parse time. We present the language and its equirecursive type
system, which is inferred with no annotations and in which the occurs check is a
diagnosis rather than a failure; we prove type soundness in a form that is
necessarily two-sided because inversion is total; and we characterise
expressiveness: the relations denotable in 42 are exactly the recursively
enumerable ones, with Turing completeness as the single-valued case.

\keywords{Reversible computation \and Relational semantics \and Point-free
programming \and Dagger categories \and Expressiveness}
\end{abstract}

\section{Introduction}

Every functional reversible language in the literature carries the same
obligation. If a program is a set of clauses and it must run backwards
deterministically, then no two clauses may produce results that overlap --- and
the languages differ mainly in how they discharge it.

\textbf{Inv}~\cite{inv} assumes it. Mu, Hu and Takeichi require that in $f \cup g$
``not only the domains, but the ranges of $f$ and $g$'' be disjoint, and add that
``the disjointness may be checked by a type system, but we have not explored this
possibility''. The condition costs them a primitive: an inequality test
\texttt{neq} exists in the language because it is ``sometimes necessary for
ensuring the disjointness of the two branches of a union''.

\textbf{Theseus}~\cite{theseus} makes it the single rule a programmer must
maintain, on both sides of every clause:

\begin{quote}\small
Non-overlapping and exhaustive coverage in pattern clauses. The collections of
patterns in the left-hand side (LHS) of each clause must be a complete
non-overlapping covering of the input type. Similarly, the collections of patterns
in the right-hand side (RHS) of each clause must also be a complete
non-overlapping covering of the return type.
\end{quote}

\textbf{RFUN}~\cite{rfun} identifies the same problem --- ``case branches may be
non-orthogonal: the result might conceivable have come from several branches'',
which it calls ``the functional variant of the problem of the general
irreversibility of if-then-else constructs'' --- and answers it not by forbidding
overlap but by ordering it, under a first-match policy in which a result ``must
never match a branch that textually precedes it''.

\textbf{PisoLang}~\cite{pisolang} decides orthogonality by unification, as two
premises of its case-typing rule. \textbf{Chardonnet, Lemonnier and
Valiron}~\cite{chardonnet} formalise it as a structural relation $v_1 \perp v_2$
appearing twice in the typing rule for an iso, keeping non-overlap while dropping
exhaustivity ``in order to allow non-terminating behaviour''. In the imperative
branch, \textbf{Janus}~\cite{janus,lutz} obliges the programmer to discharge it.
Its conditional carries a second predicate: ``the predicate after \texttt{if} is
the test, and that after \texttt{fi} is the assertion'', and ``the assertion makes
the conditional reversible'' --- at the price, stated as plainly, that ``if the
assertion does not have the required value, execution of the loop is undefined''.

Six languages, six answers, and one question none of them asks: \emph{what if the
condition is simply not imposed?}

\subsection{What happens if you decline}

Declining it means leaving the category \PInj{} of partial injections for the
category \Rel{} of sets and relations. The backward direction becomes many-valued,
and the defining law of the language acquires an $\in$:
\[
    x \in P(y) \quad\iff\quad y \in \mathrm{inv}(P)(x)
\]
Read carefully, this says that running forwards and then backwards returns a set
\emph{containing} where you started, rather than returning where you started. That
is the price, it is paid in exactly one place, and we state it at the outset
because the rest of the paper is an argument that it is worth paying.

What it buys is that inversion becomes \textbf{total}. There is no side condition,
no well-formedness check, no proof obligation and no undecidable question anywhere
in the definition of the inverse of a program. Concretely, the entire inversion
algorithm of the implementation is

\begin{lstlisting}
def dagger(t: Term) -> Term:
    match t:
        case Prim(n, i):  return Prim(n, not i)
        case Ref(n, i):   return Ref(n, not i)
        case Seq(s, u):   return Seq(dagger(u), dagger(s))   # contravariant
        case Union(s, u): return Union(dagger(s), dagger(u))
        case Sum(s, u):   return Sum(dagger(s), dagger(u))
        case Prod(s, u):  return Prod(dagger(s), dagger(u))
        case Star(s):     return Star(dagger(s))
\end{lstlisting}

\noindent and \texttt{dagger(dagger(t)) = t} holds syntactically, on the nose.
Because it is total, the surface syntax can eliminate the inversion operator
\texttt{!} at parse time, so that it never appears in the abstract syntax at all:
\[
    \texttt{parse("(copy ; join)!")} \;=\; \texttt{parse("join! ; copy!")}
\]

This property has been named in the literature at least three times, which is
itself worth recording. Theseus calls an operator \emph{syntactically reversible}
when its inverse reading coincides with its inverse meaning. Yokoyama and Gl\"uck
call an inversion \emph{local} ``iff for any given program unit the inversion can
always produce an inverse unit'', and prove Janus has it. And the 1993 design of
Section~\ref{sec:prov} states it as an equation,
$(\odot(R_1,\dots,R_n))^{-1} = \odot(R_n^{-1},\dots,R_1^{-1})$. All three give
sequential composition as the worked example, and all three arrive at the
contravariant law --- Janus as $(s_1\,s_2)\breve{\;} \sim \breve{s_2}\,\breve{s_1}$.

Where they differ is what it costs to hold the property everywhere. \emph{Janus
holds it}, and pays with the assertion: a predicate the programmer must supply for
every conditional, undefining the program when it is wrong. \emph{Theseus does not
hold it} --- its own conditional fails, and so does its boolean test, leaving open
whether restricting tests to a symmetric subclass costs expressiveness. The 1993
design does not hold it either, for the same reason.

\emph{This paper's language holds it with neither price}: nothing for the
programmer to discharge, and no construct excluded. The mechanism is that there is
no conditional to fail it and no test operator to weaken --- branching is the sum
functor, and filtering is a composite of the diagonal (Section~\ref{sec:gadgets}).

\subsection{Contributions}

\begin{enumerate}
\item \textbf{A point-free reversible language over \Rel{} in which every
  combinator is syntactically reversible} (Sections~\ref{sec:lang},
  \ref{sec:buys}). We show that this is precisely what dropping the disjointness
  condition buys, and that the resulting metatheory is smaller than that of the
  injective languages rather than larger.
\item \textbf{An equirecursive type system, inferred with no annotations}
  (Section~\ref{sec:types}), in which a cyclic substitution is admitted exactly
  when the recursive type it describes is inhabited --- the occurs check as a
  diagnosis rather than a failure. We prove type soundness, and show that the
  statement must be two-sided, because at the inversion rule the goal concerns the
  domain of the premise.
\item \textbf{A characterisation of expressiveness} (Section~\ref{sec:expr}): the
  relations denotable in 42 are exactly the recursively enumerable ones. Turing
  completeness is the single-valued case; the upper bound is the more informative
  half, and it holds because every combinator is a positive operation and nothing
  in the language forms a complement.
\end{enumerate}

\subsection{Provenance}\label{sec:prov}

The design is not new. It is the second instantiation of a language called $4_2$,
specified in a 1993 Master's thesis at the University of Amsterdam~\cite{jansen}
and never published, which arose not from reversible computing but from machine
translation --- bidirectional grammars for the Rosetta project at Philips. That
thesis already states the defining law above as a numbered equation, already
defines syntactic reversibility as the displayed equation of Section~1.1, and
already finds its own conditional violating it. What it does not have is the
point-free core: $4_2$ is imperative, with variables and assignment-shaped
transformers, and the semiring primitive set of Section~\ref{sec:lang} is the
second instantiation's contribution.

The thesis was not written in ignorance of the field it now sits in: its
bibliography cites Landauer~\cite{landauer} and Bennett~\cite{bennett} directly,
alongside the program-inversion literature and Plotkin's powerdomain
construction~\cite{plotkin} for the set-valued reading of a program. What it did
not have was a reversible-languages community to publish into; Janus existed but
was invisible until 2007. We record the lineage for honesty about priority, not
as evidence; nothing in this paper rests on it.

\section{The Language}\label{sec:lang}

\subsection{Syntax}

Types are the trees over $0$, $1$, $+$ and $\times$, with variables and a
recursion binder,
\[
    T ::= 0 \mid 1 \mid T + T \mid T \times T \mid X \mid \mu X.\,T
\]
and are \textbf{equirecursive}: $\mu X.\,F$ and $F[\mu X.F/X]$ are the same type,
not merely isomorphic. Values are finite trees over the same four constructors ---
$()$, $L\,v$, $R\,v$, $(v_1,v_2)$ --- and $V(T)$ denotes the values of type $T$. We
write $\mathbb{V}$ for the union of all of them, since the evaluator is untyped
and works on $\mathbb{V}$.

A program is a set of definitions over thirteen primitives and six combinators:
\begin{center}\small
\begin{tabular}{@{}ll@{}}
\texttt{t ; u} & sequential composition, in diagrammatic order \\
\texttt{t + u} & the sum functor \\
\texttt{t * u} & the product functor \\
\texttt{t | u} & union \\
\texttt{t\textasciicircum{}} & reflexive-transitive closure \\
\texttt{t!} & inversion \\
\end{tabular}
\end{center}
\noindent(We give the combinators as a table rather than a grammar, because the
union operator and BNF alternation are the same character.) The thirteen primitives (Fig.~\ref{fig:prims}) are the isomorphisms
witnessing the commutative-semiring structure of $(0,1,+,\times)$ --- that is, a
rig groupoid presentation --- together with the two morphisms that make the
setting \Rel{} rather than a groupoid.

\begin{figure}[t]
% \footnotesize, not \small: at \small the four columns compute to 343.5pt
% against a 346.4pt text width, which is too close to the edge to trust.
\centering\footnotesize
\begin{tabular}{@{}ll@{\qquad}ll@{}}
\hline
primitive & scheme & primitive & scheme \\
\hline
\texttt{id}        & \texttt{a <-> a}                    & \texttt{dist}   & \texttt{(a + b) x c <-> a x c + b x c} \\
\texttt{zero}      & \texttt{a <-> b}                    & \texttt{inl}    & \texttt{a <-> a + b} \\
\texttt{swapsum}   & \texttt{a + b <-> b + a}            & \texttt{inr}    & \texttt{a <-> b + a} \\
\texttt{assocsum}  & \texttt{a + (b + c) <-> a + b + c}  & \texttt{copy}   & \texttt{a <-> a x a} \\
\texttt{unitsum}   & \texttt{0 + a <-> a}                & \texttt{join}   & \texttt{a + a <-> a} \\
\texttt{swapprod}  & \texttt{a x b <-> b x a}            &                 & \\
\texttt{assocprod} & \texttt{a x (b x c) <-> a x b x c}  &                 & \\
\texttt{unitprod}  & \texttt{1 x a <-> a}                &                 & \\
\hline
\end{tabular}
\caption{The primitive schemes. This is the implementation's own table, and is
checked against it by the artifact's test suite.}
\label{fig:prims}
\end{figure}

\texttt{copy} and \texttt{join} are the diagonal and the codiagonal. Neither is an
isomorphism, and their presence is the whole of the difference between this
language and the groupoid fragment; Section~\ref{sec:related} returns to the fact
that $\Pi$ names exactly these two as identities that do \emph{not} hold, and
recovers them only as effects.

\subsection{Typing}

\begin{figure}[t]
\centering
\[
\begin{array}{c@{\qquad}c}
\rul{\vdash t : A \iso B \quad \vdash u : B \iso C}
    {\vdash t \mathbin{;} u : A \iso C} \;\textsc{seq}
&
\rul{\vdash t : A \iso B \quad \vdash u : C \iso D}
    {\vdash t + u : A + C \iso B + D} \;\textsc{sum}
\\[3.2ex]
\rul{\vdash t : A \iso B \quad \vdash u : C \iso D}
    {\vdash t * u : A \times C \iso B \times D} \;\textsc{prod}
&
\rul{\vdash t : A \iso B \quad \vdash u : A \iso B}
    {\vdash t \mid u : A \iso B} \;\textsc{union}
\\[3.2ex]
\rul{\vdash t : A \iso A}{\vdash t\clo : A \iso A} \;\textsc{star}
&
\rul{\vdash t : A \iso B}{\vdash t! : B \iso A} \;\textsc{dagger}
\end{array}
\]
\caption{Typing rules, over the primitives of Fig.~\ref{fig:prims} taken as axioms
at every ground instantiation.}
\label{fig:typing}
\end{figure}

Two remarks on Fig.~\ref{fig:typing}. There are no \texttt{fold} and
\texttt{unfold} rules, and their absence \emph{is} the equirecursive choice:
$\mu X.F$ and its unfolding are the same type, so nothing has to be written to pass
between them. And the \textsc{dagger} rule is the type-level shadow of the defining
law --- it exchanges the two sides and does nothing else, which is why it is
implemented as a swap of the inferred scheme.

\subsection{Semantics}

The evaluator is untyped: $\den{t} \subseteq \mathbb{V} \times \mathbb{V}$, and a
shape mismatch denotes the empty relation rather than raising an error. Being
outside the domain of a relation is not a failure; it is having no image.

\begin{figure}[t]
\begin{lstlisting}
[[id]]        = { (v, v) }                    [[zero]]   = {}
[[swapsum]]   = { (L v, R v) } u { (R v, L v) }
[[assocsum]]  = { (L a, L (L a)) } u { (R (L b), L (R b)) } u { (R (R c), R c) }
[[unitsum]]   = { (R v, v) }
[[swapprod]]  = { ((a, b), (b, a)) }
[[assocprod]] = { ((a, (b, c)), ((a, b), c)) }
[[unitprod]]  = { (((), v), v) }
[[dist]]      = { ((L a, c), L (a, c)) } u { ((R b, c), R (b, c)) }
[[inl]]       = { (v, L v) }                  [[inr]]    = { (v, R v) }
[[copy]]      = { (v, (v, v)) }
[[join]]      = { (L v, v) } u { (R v, v) }

[[t ; u]] = { (a, c) : exists b. (a, b) in [[t]] and (b, c) in [[u]] }
[[t + u]] = { (L a, L b) : (a, b) in [[t]] } u { (R c, R d) : (c, d) in [[u]] }
[[t * u]] = { ((a, c), (b, d)) : (a, b) in [[t]] and (c, d) in [[u]] }
[[t | u]] = [[t]] u [[u]]     [[t^]] = U_{n >= 0} [[t]]^n     [[t!]] = [[t]]^o
\end{lstlisting}
\caption{Denotations. $v$, $a$, $b$, $c$ range over $\mathbb{V}$; \texttt{u} is
union and \texttt{\^{}o} is converse.}
\label{fig:sem}
\end{figure}

$\den{\texttt{inl}}$ and $\den{\texttt{inr}}$ are total functions forwards and
partial backwards, which is where partiality enters. $\den{\texttt{join!}}$ is the
only primitive whose image can have two elements: \emph{every set larger than a
singleton anywhere in the language traces back to it}, which is worth knowing,
because it means ambiguity in 42 has exactly one origin.

A program is a finite, possibly mutually recursive system of definitions inducing
$\Phi : \Rel^n \to \Rel^n$. Every operation above is monotone and Scott-continuous
--- including converse, which is a lattice isomorphism --- so the program denotes
$\mathrm{lfp}\,\Phi = \bigcup_k \Phi^k(\emptyset,\dots,\emptyset)$.

Parameterised definitions (\texttt{def ctrl m = mat ; (id + m) ; mat!}) are
second-order: a parameter denotes a relation, never another combinator. They are
eliminated by substitution before evaluation, and we assume throughout that
programs contain none. Section~\ref{sec:related} notes that Theseus and RFUN make
the same choice and describe it in the same terms.

\subsection{This is not Kleene algebra with tests}

The operators $;$, $\mid$ and $\clo$ are the regular ones, and the resemblance
is not accidental: the 1993 design cites Pratt~\cite{pratt} and Harel~\cite{harel}
for exactly this fragment. The differences are two, and they are what the paper is
about.

First, Kleene algebra with tests is \emph{single-sorted} --- one carrier, with
tests a Boolean subalgebra of it --- whereas terms here are typed $A \iso B$
between \emph{different} objects, and $+$ and $*$ are \emph{functors} on that
typing, not operations on a single carrier. Second, the base is not an arbitrary
alphabet but a fixed, forced generating set: the semiring isomorphisms. The setting
is therefore \emph{regular control over rig-groupoid data}, and the content is in
the interaction: \texttt{dist} is what lets $;$, $\mid$ and $\clo$ see inside a
sum, and \texttt{copy}/\texttt{join} are what let them leave the groupoid.

One consequence recurs in Section~\ref{sec:expr}: over \emph{finite} types, $;$,
$\mid$ and $\clo$ give exactly finite-state power. All unboundedness comes from
the $\mu$-types, never from the control structure.

\section{What Dropping the Condition Buys}\label{sec:buys}

\subsection{The obligation, and what it costs those who accept it}

Section~1 listed six discharges of one condition. What is less often noted is what
each costs. \textbf{Inv} spends a primitive: \texttt{neq} is in the language to
establish disjointness. \textbf{Theseus} must reject programs a programmer would
expect to write; its Section~3.1 gives four invalid expressions, of which
\texttt{drop\_var} merely fails to use a bound variable on the other side and
\texttt{dup\_var} uses one twice, both forbidden by the rule that each variable
``must appear exactly once on the other side and with the same type''.
\textbf{PisoLang} must decide pattern orthogonality by unification and carve out a
controlled non-linear use of variables to recover duplication.
\textbf{Chardonnet et al.} carry orthogonality as two premises in a typing rule.
\textbf{RFUN} must fix an order on clauses and make the semantics depend on it.
\textbf{Janus} pushes the obligation onto the programmer, whose program is
undefined if the assertion is wrong.

Janus repays a closer look, because on one point it agrees with us against the
functional languages. It does \emph{not} require its operations to be injective:
``The evaluation of expressions is not backward deterministic because function
$[\![\odot]\!]$ is not injective, and thus there exists no inverse. As we shall
see, this does not harm the backward and forward determinism of Janus
statements.'' Injectivity is demanded of statements, not of arithmetic, and a
syntactic restriction on assignment buys it. So Janus and this paper agree that
non-injective operations are admissible, and differ only in where the cost is
paid: Janus in a restriction plus an assertion per conditional, 42 in the $\in$.

\subsection{What declining it buys}

Four things, and the first is the one that matters.

\textbf{Inversion is total.} \texttt{dagger} is defined on every term, by the
eight-line structural recursion of Section~1.1. There is no case in which it is
undefined, no condition it must check, and no property of the argument it must
establish first; $\texttt{dagger} \circ \texttt{dagger} = \mathrm{id}$ holds
syntactically.

\textbf{\texttt{!} need not exist in the abstract syntax.} Because inversion is
total, the parser can push it to the leaves and discard it, so no evaluator,
printer, optimiser or type rule ever encounters an inverted composite.

\textbf{There is no progress theorem to weaken.} In a language of partial
isomorphisms a stuck term is an embarrassment: PisoLang records that
\texttt{(case True <-> ()) False} is well-typed and stuck, so progress fails. Here
that term denotes $\emptyset$, a perfectly good morphism of \Rel{}, and no theorem
is lost because none was claimed.

\textbf{Ambiguity has one source.} Every non-singleton image traces to
\texttt{join!}. A language that forbids ambiguity must exclude it everywhere; one
that permits it can localise it.

\subsection{What it costs}

The $\in$. Forward-then-backward returns a set containing where you started:
\begin{center}\small
\texttt{add(2,3) = 5} \qquad
\texttt{add!(5) = \{(0,5), (1,4), (2,3), (3,2), (4,1), (5,0)\}}
\end{center}
Both directions run the \emph{same definition}; the backward one is its dagger,
computed mechanically. Whether this is acceptable is a question about what one
wants a reversible language for --- and we note that composition can restore
determinism that a factor lacks. \texttt{double = copy ; add} has a single-valued
dagger, because \texttt{copy!} keeps only pairs whose components agree, so
\texttt{double! 6 = \{3\}} and \texttt{double! 7 = \{\}}.

\subsection{Why the trade is favourable in principle}

There is a reason the metatheory gets smaller rather than larger, and it is
semantic. The class of relations 42 denotes will be shown in
Section~\ref{sec:expr} to be the recursively enumerable relations, and that class
is \textbf{closed under converse} --- swap the pairs in the enumeration. The class
an injective language denotes, computable partial injections, is closed under
inverse too, but the closure has to be \emph{earned}: an injective language must
ensure that every construct it admits preserves injectivity, which is exactly the
disjointness obligation. A language over \Rel{} gets the closure for free, because
the ambient category is already a dagger category. \emph{The totality of \texttt{!}
in the syntax is a shadow of the fact that \Rel{} has a dagger and \PInj{}'s
inverse must be constructed.}

\section{Types}\label{sec:types}

Types are inferred and never written. Nothing in the surface syntax mentions a
type, and there is no \texttt{fold}/\texttt{unfold}.

\subsection{Inference}

Algorithm W with constraints and unification. Because every definition is a closed
term, there is no environment of monomorphic assumptions to avoid capturing, and
none of the usual generalisation machinery --- levels, ranks, the value restriction
--- is needed. Mutually recursive groups are solved as a unit against monomorphic
assumptions and generalised once, Milner-style.

A $\mu$ is never something a programmer reaches for. It is what the checker infers
when unification closes a loop: \texttt{def double = copy ; add} is inferred at
$\mu X.\,1 + X \iso \mu Y.\,1 + Y$. Here \texttt{nat} is declared nowhere;
\texttt{copy} forces \texttt{add}'s two summands to agree, and that \emph{is}
\texttt{nat <-> nat}.

\subsection{A cycle is the recursive type}

Where a finite-tree checker runs an occurs check and gives up, this one binds the
variable to a term containing itself and carries on. Unification then works up to
unfolding, \emph{coinductively}: assume the two sides equal, and check that the
assumption survives one layer. The set of assumptions is finite because the regular
trees a cyclic substitution can describe have finitely many distinct subterms, and
that is what makes the procedure terminate.

\subsection{But a cycle is read before it is trusted}

Not every cycle is a type. $b = a + b$ is $\mu X.\,a + X$, a list, which has finite
values; $a = a \times a$ is $\mu X.\,X \times X$, which has none --- and
\texttt{copy\textasciicircum{}} is how one writes it. The test is one unfolding with the variable
taken empty: \emph{$\mu X.\,F(X)$ has a finite value exactly when $F(0)$ does.} If
$F(0)$ is empty then so is $F(F(0))$, and so on. The occurs check therefore becomes
a \emph{diagnosis}: a cycle is admitted when the recursive type it describes is
inhabited and rejected when it is not, and the rejection names the equation that
could not be satisfied. We know of no other language in this family that infers its
recursive types at all; all of them are declared.

\subsection{Soundness, and why it must be two-sided}

\begin{proposition}[Soundness]\label{prop:sound}
Let $\vdash t : A \iso B$ with $A$, $B$ ground. Then for every
$(v,w) \in \den{t}$,
\[
    v \in V(A) \;\text{ or }\; w \in V(B)
    \quad\text{implies}\quad
    v \in V(A) \;\text{ and }\; w \in V(B).
\]
\end{proposition}

The expected form --- \emph{$v \in V(A)$ implies $w \in V(B)$} --- \textbf{cannot be
proved by induction here}. At the \textsc{dagger} rule the goal for $t!$ concerns
the codomain of $t!$, which is the \emph{domain} of $t$, and the induction
hypothesis says nothing about domains. The statement must therefore carry both
directions at once. Stated as above it is \textbf{self-dual}: the claim for
$t : A \iso B$ and the claim for $t! : B \iso A$ are literally the same sentence,
so the dagger case is discharged by observing that $\den{t!} = \den{t}\conv$ and
that the condition is symmetric. This is the defining law appearing in the
metatheory: \emph{a language whose inversion is total cannot have a one-sided type
soundness theorem, and does not need one.}

\begin{proof}
Induction on the derivation. The thirteen primitive cases are immediate from
Fig.~\ref{fig:sem}, and three are representative. For $\texttt{join} : A' + A' \iso
A'$ a pair is $(L\,a, a)$ or $(R\,a, a)$, and $L\,a \in V(A'+A') \iff a \in V(A')$
--- the biconditional the statement wants. For $\texttt{inl} : A \iso A + B'$ a pair
is $(v, L\,v)$ and $L\,v \in V(A+B') \iff v \in V(A)$. For
$\texttt{unitsum} : 0 + A \iso A$ a pair is $(R\,v, v)$, and $V(0+A)$ has no $L$
component at all because $V(0) = \emptyset$ --- which is exactly why the evaluator
returns the empty set on an $L$. \texttt{zero} is vacuous.

For \textsc{seq}, let $(a,c) \in \den{t\mathbin{;}u}$, so $(a,x) \in \den{t}$ and
$(x,c) \in \den{u}$. If $a \in V(A)$ the hypothesis for $t$ gives $x \in V(B)$ and
then that for $u$ gives $c \in V(C)$; if instead $c \in V(C)$, the hypotheses are
used in the other order. Both readings are available precisely because the
statement is two-sided. \textsc{sum} and \textsc{prod} reduce to the hypotheses
componentwise, and \textsc{prod} needs the two-sided form even with no dagger in
sight, since knowing only $(b,d) \in V(B \times D)$ requires reading each
hypothesis backwards. \textsc{union} applies either hypothesis. For \textsc{star},
$A = B$; $n = 0$ is the identity and each further step is the hypothesis for $t$ in
whichever direction the given end supplies. \textsc{dagger} is as above. For a
recursive group, $\den{f_i} = \bigcup_k \Phi^k(\emptyset,\dots,\emptyset)_i$; each approximant
satisfies the property, and the property is preserved by unions of chains, being a
condition on individual pairs. \qed
\end{proof}

\begin{corollary}\label{cor:cod}
If $\vdash t : A \iso B$ and $v \in V(A)$ then every result of running $t$ forwards
on $v$ lies in $V(B)$; and if $w \in V(B)$, every result of running it backwards
lies in $V(A)$.
\end{corollary}

Note what this does \emph{not} say. Nothing here promises a result exists:
$\den{\texttt{pred}}$ at $0$ is empty, which is a legitimate morphism, so there is
no progress property to state and none is lost.

\section{Expressiveness}\label{sec:expr}

\begin{theorem}\label{thm:main}
Let $A$, $B$ be closed types and $R \subseteq V(A) \times V(B)$. There is a closed
42 program denoting $R$ \textbf{iff} $R$ is recursively enumerable.
\end{theorem}

Two remarks on why this and not ``42 is Turing complete'', which is the
$\supseteq$ half alone. First, the $\subseteq$ half is the more informative one: it
says 42 denotes \emph{exactly} $\Sigma^0_1$, so no program decides the complement
of the halting problem and no extension of the primitive table can change that
without adding a form of negation. Second, the statement has the right shape for a
relational language. The reference result is Axelsen and Gl\"uck's~\cite{axelsen},
that reversible Turing machines compute exactly the injective computable functions;
Chardonnet et al.\ prove the corresponding statement for a Theseus-descended
language, concluding that it ``fully characterises all of the computable morphisms
in PInj''. Read together, their theorem and Theorem~\ref{thm:main} are the same
statement in two different categories --- \emph{every computable morphism of the
ambient category is denotable} --- with \PInj{} there and \Rel{} here.
\textbf{What 42 changes is not the theorem; it is the category.}

\subsection{Soundness}

\begin{proposition}\label{prop:re}
For every program and every closed term $t : A \iso B$ over it, $\den{t}$ is
recursively enumerable, and an index is computable from the text.
\end{proposition}

\begin{proof}[sketch]
Three steps. The primitives denote decidable relations on a decidable set of finite
trees. The six operations preserve r.e.\ uniformly: composition by dovetailing,
union by interleaving, $+$ and $*$ by applying the tag or the pairing, converse by
swapping each enumerated pair, and $\den{t\clo}$ because $\den{t}^n$ is
uniformly r.e.\ in $n$. Recursion adds nothing: by continuity
$\mathrm{lfp}\,\Phi = \bigcup_k \Phi^k(\emptyset,\dots,\emptyset)$, an index for each
approximant is computable from $k$, and dovetailing gives the limit. \qed
\end{proof}

\begin{corollary}
No 42 program denotes a relation that is not $\Sigma^0_1$.
\end{corollary}

\textbf{Where the bound comes from} is worth naming, because it is not where one
might expect. Not from reversibility, and not from any restriction on recursion ---
42 has general recursion. It comes from the shape of the syntax: every combinator
is a \emph{positive} operation, and \textbf{nothing in the language forms a
complement}. A hypothetical combinator ``$t$ fails here'' would denote the
complement of a domain and break Proposition~\ref{prop:re} at once. The practical
consequence is that the primitive table may be enlarged with any decidable relation
and the theorem survives.

It is a small pleasure that the 1993 design excluded complement too, for a related
but distinct reason: that the complement of a finite result set in an infinite
universe is not finitely presentable. Convergent instinct, different argument.

\subsection{Four definable gadgets}\label{sec:gadgets}

\begin{lemma}[Discard]\label{lem:drop}
For every closed $C$ there is a term $\mathit{drop}_C : C \iso 1$ with
$\den{\mathit{drop}_C} = V(C) \times \{()\}$.
\end{lemma}

by $\mathit{drop}_0 = \texttt{zero}$, $\mathit{drop}_1 = \texttt{id}$,
$\mathit{drop}_{A+B} = (\mathit{drop}_A + \mathit{drop}_B)\mathbin{;}\texttt{join}$,
$\mathit{drop}_{A\times B} = (\mathit{drop}_A * \mathit{drop}_B)\mathbin{;}\texttt{unitprod}$,
and a recursive definition following $F$ for $\mu X.F$; totality on the last is an
induction on the size of the finite tree.

This deserves a comment. \texttt{discard} is the operation the language is designed
to avoid --- the rule is that nothing may be thrown away --- and there is no such
primitive. It is definable anyway, and there is no contradiction: the rule is that
if a step loses information then the backward direction must be allowed to return
several candidates, and $\mathit{drop}!$ returns all of them. \emph{42 forbids
forgetting reversibly, not forgetting.} By contrast Theseus lists precisely this
program, \texttt{drop\_var}, among its invalid expressions.

\begin{lemma}[Universal relation]\label{lem:univ}
$\mathit{univ} = \mathit{drop}_A \mathbin{;} \mathit{drop}_B!$ denotes
$V(A) \times V(B)$.
\end{lemma}

\begin{lemma}[Filter]\label{lem:filt}
For any $\mathit{test} : C \iso 1$ with domain $U$, the term
$\texttt{copy} \mathbin{;} (\mathit{test} * \texttt{id}) \mathbin{;} \texttt{unitprod}$
denotes the partial identity on $U$.
\end{lemma}

This is how a semi-decision becomes a program, and it is also the answer to
Theseus's open question of Section~1.1: the filter is not a construct but a
composite, it is its own dagger for \emph{any} test, and no restriction to a
symmetric subclass is needed.

\begin{lemma}[Serialisation]\label{lem:ser}
For every closed $C$ there is $\mathit{ser}_C : C \iso \mathit{bits}$ denoting the
graph of an injective computable $e_C : V(C) \to \{0,1\}^*$ whose image is
prefix-free and decidable.
\end{lemma}

by $\mathit{ser}_1 = \texttt{inl}$,
$\mathit{ser}_{A+B} = ((\mathit{ser}_A\mathbin{;}\mathit{cons}_0) + (\mathit{ser}_B\mathbin{;}\mathit{cons}_1))\mathbin{;}\texttt{join}$,
$\mathit{ser}_{A\times B} = (\mathit{ser}_A * \mathit{ser}_B)\mathbin{;}\mathit{append}$,
and recursion for $\mu$. Prefix-freeness is the same induction: two prefixes of one
string are comparable, so concatenation of prefix-free codes is prefix-free.

The product clause is worth pausing on because it looks unsound and is not.
$\mathit{append}$ run backwards yields \emph{every} way to split a list --- six, for
the five-bit string encoding $(1,2)$ --- and yet $\mathit{ser}_{nat \times nat}$ run
backwards yields exactly one, because prefix-freeness leaves only the split that
decodes. Composition restoring determinism that a factor lacks is here doing real
work in a proof.

\subsection{Simulation}

\begin{lemma}[Halting]\label{lem:halt}
For every deterministic Turing machine $M$ there is a term
$\mathit{halts}_M : \mathit{bits} \iso 1$ with
$\den{\mathit{halts}_M} = \{(x,()) : M \text{ halts on } x\}$.
\end{lemma}

We assume $M$ normalised to a binary alphabet, a one-way infinite tape it never
falls off, and a single halting state; all three are standard~\cite{hopcroft} and
preserve the halting set.

\paragraph{The encoding.} A tape is
$\mathit{bits} \times ((1+1) \times \mathit{bits})$ --- left reversed, head, right
--- and a configuration is
$\mathit{conf} = \mathit{tape} + (\mathit{tape} + (\cdots + \mathit{tape}))$ with
one summand per state. The control state is a \textbf{label, not a component of a
pair}: to act differently in each state one then writes $f + g + h$, which is the
sum functor doing exactly what it is for, and no plumbing is needed to reach the
state and put it back. This one decision is what keeps the transition relation to a
line.

Head movement pushes the head cell onto the left list and pops the right, and
reading the head brings it to the front and splits on it:

\begin{lstlisting}
right    = assocprod ; ((swapprod ; inr) * id) ; (id * inr!)
readhead = focus ; dist ; (unitprod + unitprod)
\end{lstlisting}

\noindent\textbf{\texttt{left} is not written out --- it is \texttt{right!}}, and
writing the head is \texttt{readhead} run backwards. For
$\delta(q,s) = (q',s',m)$:

\begin{lstlisting}
beta_{q,s} = write_{s'} ; move_m ; tag_{q'}
delta_q    = readhead ; (beta_{q,0} + beta_{q,1}) ; join     delta_{q_h} = zero
step       = (delta_{q_1} + (delta_{q_2} + ...)) ; collapse_h
halts_M    = load ; tag_{q_1} ; step^ ; tag_{q_h}! ; drop_tape
\end{lstlisting}

\paragraph{Correctness} is four lemmas. $\den{\mathit{step}}$ is a partial function;
it simulates $\delta$ on representations; the halting state is a dead end because
$\delta_{q_h} = \texttt{zero}$; and therefore $\den{\mathit{step}\clo}$ is
exactly the run, from which $\mathit{tag}_{q_h}!$ selects the halted configuration,
of which minimality of the halting time gives at most one.

The one case worth keeping in full is where the finite/infinite mismatch is paid
for. A Turing tape is infinite and a value is a finite tree, so the encoding is
\textbf{many-to-one}: trailing blanks are represented by their absence. When the
head moves right off the end of the finite list, a padding branch invents a blank
--- and this is correct precisely because an empty right-list forces the next cell
to be blank. Note what the simulation lemma does \emph{not} claim: it says nothing
about the successor representation being canonical. Feed it a representation with
three trailing blanks and it returns one with three trailing blanks, and since the
statement is about the \emph{reading}, the padding simply rides along. The mismatch
costs one branch of one definition and one case of one proof.

An instance of this construction --- a three-state machine incrementing a binary
number --- is in the artifact and runs.

\subsection{Completeness}

\begin{theorem}\label{thm:complete}
For closed $A$, $B$ and r.e.\ $R \subseteq V(A) \times V(B)$, some closed 42
program denotes $R$.
\end{theorem}

\begin{proof}
Put $C = A \times B$, so $R \subseteq V(C)$. By Lemma~\ref{lem:ser},
$S = e_C(R) \subseteq \{0,1\}^*$ is r.e., since $e_C$ is a computable injection with
decidable image. Take $M$ halting exactly on $S$; by Lemma~\ref{lem:halt},
$\mathit{test} = \mathit{ser}_C \mathbin{;} \mathit{halts}_M$ has domain $R$. By
Lemma~\ref{lem:filt},
$\mathit{filt} = \texttt{copy} \mathbin{;} (\mathit{test} * \texttt{id}) \mathbin{;} \texttt{unitprod}$
is the partial identity on $R$. With $\mathit{univ}$ from Lemma~\ref{lem:univ}, put
\[
    t \;=\; \texttt{copy} \mathbin{;} (\texttt{id} * \mathit{univ}) \mathbin{;}
            \mathit{filt} \mathbin{;} (\mathit{drop}_A * \texttt{id}) \mathbin{;}
            \texttt{unitprod}
\]
and trace $x \in V(A)$: \texttt{copy} gives $(x,x)$; $(\texttt{id} * \mathit{univ})$
gives $(x,y)$ for \emph{every} $y \in V(B)$; $\mathit{filt}$ keeps those with
$(x,y) \in R$; $(\mathit{drop}_A * \texttt{id})$ gives $((),y)$; \texttt{unitprod}
gives $y$. So $\den{t} = R$. \qed
\end{proof}

The shape of that argument is the point: \textbf{guess the output, then check it.}
The guess is $\mathit{drop}_A \mathbin{;} \mathit{drop}_B!$ and the check is
\texttt{copy} --- both already in the language, neither invented for the proof. An
injective language cannot write the guess at all.

\begin{corollary}[Turing completeness]
Every partial computable $f : V(A) \rightharpoonup V(B)$ is denoted by some program,
its graph being r.e.; and by Proposition~\ref{prop:re} no program denotes anything
worse.
\end{corollary}

\section{Related Work}\label{sec:related}

\paragraph{The axis.} One dial --- how many answers running backwards may give ---
orders the field.

\begin{center}\small
\begin{tabular}{@{}llp{5.8cm}@{}}
\hline
setting & $|P(y)|$ & languages \\
\hline
Groupoid, total bijections & exactly 1 & $\Pi$ \\
\PInj{}, partial injections & $\le 1$ & Janus, Theseus, RFUN, Inv,
                                        Chardonnet et al., PisoLang \\
\Rel{}, relations & arbitrary & \textbf{this paper} \\
\hline
\end{tabular}
\end{center}

Almost the whole field is the middle row: reversible languages are, with near
unanimity, \emph{injective} languages, differing in surface syntax and in how
injectivity is enforced rather than in what they denote.

\paragraph{Inv} is the nearest relative and deserves generosity. Mu, Hu and
Takeichi's point-free relational setting is this one; their converse law is this
law; and their observation that undoing a duplication requires an equality test ---
``\texttt{dup} takes a pair and lets it go through only if the two components are
equal'' --- is the observation this paper makes about \texttt{copy!}, twenty years
earlier. What differs is that Inv restricts to injective functions and 42 does not.

\paragraph{The $\Pi$ branch.} $\Pi$'s terms~\cite{pi} are the isomorphisms
witnessing that ``the structure $(b,+,0,*,1)$ is a commutative semiring'', which is
Fig.~\ref{fig:prims} minus two rows. The two rows are the interesting ones. Having
listed the isomorphisms, \emph{Information Effects} names two identities that are
\emph{not} among them --- $b \times b \not\iso b$ and
$b_1 + (b_2 \times b_3) \not\iso (b_1+b_2) \times (b_1+b_3)$ --- and 42's
\texttt{copy} and \texttt{join} are witnesses for the first and its dual. $\Pi$
does recover both, but only in an arrow metalanguage, as \emph{information effects}
built from \texttt{create} and \texttt{erase}; \texttt{join} there ``converts the
input $b+b$ to $(1+1) \times b$ and then erases the first component''. So:
\textbf{what $\Pi$ must treat as an effect, requiring a type-and-effect system and
a metalanguage, this language has as a primitive} --- because in \Rel{} those two
morphisms are morphisms. Two further differences: $\Pi$ is isorecursive and pays
\texttt{fold}/\texttt{unfold}, where equirecursive types pay nothing; and $\Pi$'s
$\mathit{distrib0} : 0 * b \iso 0$ is absent here, $0 \times a$ being uninhabited
and so extensionally \texttt{zero}.

\paragraph{Second-order parameters.} Theseus's parametrised maps are ``a macro or a
meta-language construction'' --- ``Theseus does not have high-order maps in the
formal sense'' --- and RFUN's functional parameters are ``not enough to make
functions first class citizens''. This paper's parameters are the same choice, made
independently and for the same reason. PisoLang and Chardonnet et al.\ nest arrows
freely, the latter listing higher-order among its contributions.

\paragraph{Janus} is the imperative ancestor and the only prior language we know of
that holds local invertibility for every construct. It buys it with the assertion;
we buy it by admitting many-valued backward results. The two prices are the
alternatives this paper is really comparing, and Janus's is the more familiar one:
a proof obligation, discharged per program, whose failure mode is an undefined
execution rather than a rejected one.

\paragraph{Program inversion} is the classical ancestor: Dijkstra's
\texttt{EWD671}~\cite{ewd671}, Gries's \emph{Science of Programming}~\cite{gries},
Gries and van de Snepscheut on inverting an inorder traversal~\cite{griessnep}, and
Chen and Udding~\cite{chen}. Where that tradition inverts a \emph{given} program by
hand, 42 makes the inverse fall out of the syntax.

\section{Conclusion}

We have presented a reversible language that declines the injectivity its
neighbours buy, and argued that the trade is favourable: inversion becomes total,
\texttt{!} leaves the abstract syntax, no progress theorem is at risk, ambiguity
has a single source, and the price is one $\in$. We have given an equirecursive
type system inferred without annotation, proved it sound in the two-sided form that
a total inversion forces, and characterised the denotable relations as exactly the
recursively enumerable ones --- the same theorem the injective languages prove, read
in a larger category.

\paragraph{Limitations}, stated here rather than left to reviewers. Nothing is
mechanised; no proof assistant has seen these proofs. Three standard Turing machine
normalisations are cited, not reproved. Closure is computed by naive saturation, so
it terminates only on finite orbits: a relation can be computable in one direction
and not the other, and \texttt{pred\textasciicircum{}} saturates where \texttt{succ\textasciicircum{}} does not.
Backward runs can be exponential; nothing prunes the search. Recursion is
monomorphic, Milner-style. Inferred types are not canonical, equirecursive types
having no unique syntactic form. And parameterised definitions are eliminated by
substitution, so a \emph{recursive} combinator cannot be run even when it
typechecks: \texttt{def map f = id + (f * map f)} is well-typed at
$(a \iso b) \to (\mathit{list}\,a \iso \mathit{list}\,b)$ and its dagger is computed
correctly, but it has no finite expansion. Interpreting application rather than
expanding it away would close this, and nothing in the type system stands in the
way.

\paragraph{Future work.} The last limitation is the most inviting. Beyond it: the
denotable class is $\Sigma^0_1$ precisely because nothing forms a complement, which
suggests a principled account of what a complement-forming extension would cost.

\paragraph{Q42.} \Rel{} is matrices over the Boolean semiring. Replacing
$\mathbb{B}$ by $\mathbb{C}$ turns the same terms into finite-dimensional
unitaries, and the implementation shares the values, terms, parser, type engine and
inversion function itself with the classical one. Two constructors must go, union
and closure, being exactly those whose meaning needs addition to be
idempotent~\cite{sqrtpi}. That is a separate paper.

\paragraph{Artifact.} An implementation with no dependencies, 363 tests, every
claim in the manual re-run by a checker, and the two figures of
Section~\ref{sec:lang} verified against the implementation by the test suite.

\bibliographystyle{splncs04}
\bibliography{refs}

\end{document}
